\Askhsh{21193_SOLUTION}{1}
\begin{ylh}{1}
    \item [Λύση]
\end{ylh}
\begin{alist}
\item
    \begin{rlist}
    \item Επειδή ο ιμάντας εφάπτεται στους κυκλικούς τροχούς, οι ακτίνες $\mathrm{A\Lambda}$ και $\mathrm{\Gamma M}$ είναι ίσες και παράλληλες αφού είναι κάθετες στο ίδιο εφαπτόμενο τμήμα $\mathrm{\Lambda M}$, συνεπώς το τετράπλευρο $\mathrm{A\Lambda M\Gamma}$ είναι ορθογώνιο.
    \item Για τις γωνίες με κορυφή το κέντρο $\mathrm{A}$ του ενός τροχού έχουμε:
    \[ \mathrm{K\widehat{A}L} + 90^{\circ} + \mathrm{\widehat{A}} + 90^{\circ} = 360^{\circ} \text{ ή } \mathrm{K\widehat{A}L} + \mathrm{\widehat{A}} = 180^{\circ} \]
    δηλαδή η γωνία $\mathrm{K\widehat{A}L}$ και η γωνία $\mathrm{\widehat{A}}$ του τριγώνου $\mathrm{AB\Gamma}$ είναι παραπληρωματικές.
    \end{rlist}
\item Με βάση το α)ii. ερώτημα έχουμε:
    \[ \mathrm{\widehat{\omega}} + \mathrm{\widehat{A}} = 180^{\circ} \]
    Ανάλογα για τις γωνίες $\mathrm{\widehat{B}}$ και $\mathrm{\widehat{\Gamma}}$ βρίσκουμε $\mathrm{\widehat{\theta}} + \mathrm{\widehat{B}} = 180^{\circ}$ και $\mathrm{\widehat{\varphi}} + \mathrm{\widehat{\Gamma}} = 180^{\circ}$.
    Προσθέτοντας τις παραπάνω ισότητες προκύπτει ότι:
    \begin{align*}
        \mathrm{\widehat{\omega}} + \mathrm{\widehat{A}} + \mathrm{\widehat{\theta}} + \mathrm{\widehat{B}} + \mathrm{\widehat{\varphi}} + \mathrm{\widehat{\Gamma}} &= 180^{\circ} + 180^{\circ} + 180^{\circ} \text{ ή} \\
        \mathrm{\widehat{\omega}} + \mathrm{\widehat{\theta}} + \mathrm{\widehat{\varphi}} + 180^{\circ} &= 180^{\circ} + 180^{\circ} + 180^{\circ} \text{ ή} \\
        \mathrm{\widehat{\omega}} + \mathrm{\widehat{\theta}} + \mathrm{\widehat{\varphi}} &= 360^{\circ}
    \end{align*}
\item Με ανάλογο τρόπο, όπως για το τετράπλευρο $\mathrm{A\Lambda M\Gamma}$ του α) i. ερωτήματος, αποδεικνύεται ότι και τα τετράπλευρα $\mathrm{\Gamma NPB}$ και $\mathrm{B\Sigma KA}$ είναι επίσης ορθογώνια, οπότε θα έχουν τις απέναντι πλευρές τους ίσες, δηλαδή $\mathrm{NP} = \mathrm{\alpha}$, $\mathrm{\Lambda M} = \mathrm{\beta}$ και $\mathrm{K\Sigma} = \mathrm{\gamma}$. Αν συμβολίσουμε με $\ell_{\overset{\frown}{\mathrm{K\Lambda}}}$, $\ell_{\overset{\frown}{\mathrm{MN}}}$ και $\ell_{\overset{\frown}{\mathrm{P\Sigma}}}$ τα μήκη των μικρότερων του ημικυκλίου τόξων $\overset{\frown}{\mathrm{K\Lambda}}$, $\overset{\frown}{\mathrm{MN}}$ και $\overset{\frown}{\mathrm{P\Sigma}}$ αντίστοιχα, τότε το μήκος $\mathrm{L}$ του ιμάντα είναι:
    \begin{align*}
        \mathrm{L} &= \ell_{\overset{\frown}{\mathrm{K\Lambda}}} + \mathrm{\Lambda M} + \ell_{\overset{\frown}{\mathrm{MN}}} + \mathrm{NP} + \ell_{\overset{\frown}{\mathrm{P\Sigma}}} + \mathrm{K\Sigma} = \\
        &= \ell_{\overset{\frown}{\mathrm{K\Lambda}}} + \ell_{\overset{\frown}{\mathrm{MN}}} + \ell_{\overset{\frown}{\mathrm{P\Sigma}}} + (\mathrm{\alpha} + \mathrm{\beta} + \mathrm{\gamma}) = \\
        &= \pi \mathrm{R} \frac{\mathrm{\widehat{\omega}}}{180^{\circ}} + \pi \mathrm{R} \frac{\mathrm{\widehat{\varphi}}}{180^{\circ}} + \pi \mathrm{R} \frac{\mathrm{\widehat{\theta}}}{180^{\circ}} + 2\mathrm{\tau} = \\
        &= \pi \mathrm{R} \cdot \frac{1}{180^{\circ}} \cdot (\mathrm{\widehat{\omega}} + \mathrm{\widehat{\varphi}} + \mathrm{\widehat{\theta}}) + 2\mathrm{\tau} = \\
        &= \pi \mathrm{R} \cdot \frac{1}{180^{\circ}} \cdot 360^{\circ} + 2\mathrm{\tau} = 2\pi \mathrm{R} + 2\mathrm{\tau} = 2(\mathrm{\tau} + \pi \mathrm{R})
    \end{align*}
\end{alist}